\documentclass[a4paper,11pt]{article}

\usepackage{carnotex}

\title{Test carnotex}
\date{\today}

\begin{document}
%%%%%%%%%%%%%%%%%%%%%%%%%%%%%%%%%%%%%%%%%%%%%%%%%%%%%%%%%%%%%%%%%%%%%%%%%%%%%%
% Test de validation de carnotex.sty
%
% Ce document doit exercer toutes les commandes publiques du package
% carnotex et uniquement celles-ci.
%%%%%%%%%%%%%%%%%%%%%%%%%%%%%%%%%%%%%%%%%%%%%%%%%%%%%%%%%%%%%%%%%%%%%%%%%%%%%%
\maketitle

\section{Cadres}

\subsection{Version minimale}

\cadre{CarnoDef}{
Cadre minimal.
}

\subsection{Version complète}

\cadre{CarnoDef}[Titre]{
Cadre avec titre.
}

\subsection{Cadres imbriqués}

\cadre{CarnoProp}{
Texte.

\cadre{CarnoRem}{
Cadre interne.
}
}

\section{Lignes}

\colorline

\colorline[CarnoProp]

\colorline*[CarnoDef][6cm][1pt]

\section{Historique}

\histo{
Texte historique.
}

\section{Exercices}

\exobandeau{Exercice simple}{2}

Énoncé.

\exobandeau[5]{Exercice noté}{4}

Énoncé.

\section{Colonnes}

\col{
Colonne gauche.
}{
Colonne droite.
}

\col[0.3][0.65]{
Colonne gauche réduite.
}{
Colonne droite large.
}

\section{Utilitaires}

\siecle{21}

\siecle*{3}

\section{Bordure}

\bordure[CarnoThm]{Test}

\end{document}